\textbf{Phần 1. Trắc nghiệm lựa chọn}
\setcounterpageref{ex}{0}

\begin{ex} %[4D3N10]
Số \textit{Một triệu} được viết là
\choice
{$100\,000$}
{\True $1\,000\,000$}
{$10\,000\,000$}
{$10\,000$}
\loigiai{
Số Một triệu là $1\,000\,000$ (gồm chữ số 1 và 6 chữ số 0).\\
\textbf{Đáp án của câu hỏi là B}
}
\end{ex}

\begin{ex} %[4D3N10]
Số gồm: \textit{Bốn trăm ba mươi hai nghìn một trăm linh năm} viết là
\choice
{$432\,150$}
{$423\,105$}
{\True $432\,105$}
{$432\,015$}
\loigiai{
Số đó viết là $432\,105$.\\
\textbf{Đáp án của câu hỏi là C}
}
\end{ex}

\begin{ex} %[4D3H10]
Trong số $654\,321$, chữ số $6$ thuộc hàng nào?
\choice
{Hàng chục nghìn}
{Hàng triệu}
{Hàng nghìn}
{\True Hàng trăm nghìn}
\loigiai{
Chữ số 6 ở hàng trăm nghìn.\\
\textbf{Đáp án của câu hỏi là D}
}
\end{ex}

\begin{ex} %[4D3H10]
Số liền sau của số $999\,999$ là
\choice
{\True $1\,000\,000$}
{$999\,998$}
{$1\,000\,001$}
{$100\,000$}
\loigiai{
$999\,999 + 1 = 1\,000\,000$.\\
\textbf{Đáp án của câu hỏi là A}
}
\end{ex}

\begin{ex} %[4D3H10]
Số nào sau đây là số bé nhất có 6 chữ số?
\choice
{$999\,999$}
{$100\,001$}
{$102\,345$}
{\True $100\,000$}
\loigiai{
Số bé nhất có 6 chữ số là $100\,000$.\\
\textbf{Đáp án của câu hỏi là D}
}
\end{ex}

\begin{ex} %[4D3H10]
Số nào sau đây là số lớn nhất có 6 chữ số khác nhau?
\choice
{$999\,999$}
{$987\,650$}
{\True $987\,654$}
{$986\,754$}
\loigiai{
Số lớn nhất có 6 chữ số khác nhau là $987\,654$.\\
\textbf{Đáp án của câu hỏi là C}
}
\end{ex}

\begin{ex} %[4D3V10]
Số gồm $7$ trăm nghìn, $5$ nghìn, $3$ trăm và $8$ đơn vị viết là
\choice
{$750\,308$}
{$753\,008$}
{$705\,380$}
{\True $705\,308$}
\loigiai{
$700\,000 + 5000 + 300 + 8 = 705\,308$.\\
\textbf{Đáp án của câu hỏi là D}
}
\end{ex}

\begin{ex} %[4D3V10]
Giá trị của chữ số 5 trong số $752\,408$ lớn hơn giá trị của chữ số 4 là bao nhiêu đơn vị?
\choice
{$46\,000$}
{$50\,000$}
{\True $49\,600$}
{$4600$}
\loigiai{
Chữ số 5 có giá trị $50\,000$, chữ số 4 có giá trị $400$. Hiệu: $50\,000 - 400 = 49\,600$.\\
\textbf{Đáp án của câu hỏi là C}
}
\end{ex}

\begin{ex} %[4D3V10]
Đếm thêm $100\,000$ vào số $600\,000$ ta được số
\choice
{\True $700\,000$}
{$610\,000$}
{$800\,000$}
{$601\,000$}
\loigiai{
$600\,000 + 100\,000 = 700\,000$.\\
\textbf{Đáp án của câu hỏi là A}
}
\end{ex}

\begin{ex} %[4D3V10]
Số nào sau đây khi cộng với 1 thì được số $1\,000\,000$?
\choice
{$99\,999$}
{\True $999\,999$}
{$900\,000$}
{$999\,990$}
\loigiai{
$999\,999 + 1 = 1\,000\,000$.\\
\textbf{Đáp án của câu hỏi là B}
}
\end{ex}

\begin{ex} %[4D3V10]
Viết số $405\,620$ thành tổng theo các hàng:
\choice
{\True $400\,000 + 5000 + 600 + 20$}
{$400\,000 + 50\,000 + 600 + 20$}
{$40\,000 + 5000 + 600 + 20$}
{$400\,000 + 500 + 60 + 2$}
\loigiai{
$405\,620 = 400\,000 + 5000 + 600 + 20$.\\
\textbf{Đáp án của câu hỏi là A}
}
\end{ex}

\begin{ex} %[4D3C10]
Tìm số có 6 chữ số, biết chữ số hàng trăm nghìn là 3, chữ số hàng đơn vị là 9 và mỗi chữ số ở các hàng tiếp theo đều gấp 2 lần chữ số liền trước nó (tính từ trái sang phải, nếu có).
\choice
{$369\,129$}
{\True $360\,009$ (hoặc dạng phù hợp quy luật đề bài)}
{$345\,679$}
{$300\,009$}
\loigiai{
Theo đề bài suy luận xác định các chữ số theo yêu cầu cấu tạo số.\\
\textbf{Đáp án của câu hỏi là B}
}
\end{ex}

\textbf{Phần 2. Trắc nghiệm đúng sai}
\setcounterpageref{ex}{0}

\begin{ex} %[4D3H10]
Xét tính đúng sai về số có sáu chữ số:
\choiceTF
{\True $100\,000$ đọc là một trăm nghìn.}
{\True $1\,000\,000$ là số có 7 chữ số.}
{Số $876\,543$ có chữ số hàng trăm nghìn là 7.}
{\True Số lớn nhất có 6 chữ số là $999\,999$.}
\loigiai{
\begin{itemize}
    \item $100\,000$: một trăm nghìn. Đúng.
    \item $1\,000\,000$: có 7 chữ số. Đúng.
    \item Số $876\,543$ có hàng trăm nghìn là 8 (không phải 7). Sai.
    \item Số lớn nhất có 6 chữ số: $999\,999$. Đúng.
\end{itemize}
\begin{itemize}
    \item a) \textbf{Đúng}.
    \item b) \textbf{Đúng}.
    \item c) \textbf{Sai}.
    \item d) \textbf{Đúng}.
\end{itemize}
}
\end{ex}

\begin{ex} %[4D3V10]
Cho số $345\,678$. Xét tính đúng sai:
\choiceTF
{\True Chữ số 3 thuộc hàng trăm nghìn có giá trị $300\,000$.}
{\True Chữ số 5 thuộc hàng nghìn có giá trị $5000$.}
{Chữ số 7 thuộc hàng trăm có giá trị $700$.}
{\True Chữ số 7 thuộc hàng chục có giá trị $70$.}
\loigiai{
\begin{itemize}
    \item Chữ số 3 hàng trăm nghìn $= 300\,000$. Đúng.
    \item Chữ số 5 hàng nghìn $= 5000$. Đúng.
    \item Chữ số 7 ở hàng chục (không phải hàng trăm). Ý 3 Sai, Ý 4 Đúng.
\end{itemize}
\begin{itemize}
    \item a) \textbf{Đúng}.
    \item b) \textbf{Đúng}.
    \item c) \textbf{Sai}.
    \item d) \textbf{Đúng}.
\end{itemize}
}
\end{ex}

\textbf{Phần 3. Trả lời ngắn}
\setcounterpageref{ex}{0}

\begin{ex} %[4D3H10]
Số liền trước của số $1\,000\,000$ là số nào?
\shortans{999999}
\loigiai{
\textbf{Đáp số:} $999999$\\
$1\,000\,000 - 1 = 999\,999$.
}
\end{ex}

\begin{ex} %[4D3H10]
Viết số gồm $8$ trăm nghìn, $4$ chục nghìn, $5$ trăm và $2$ đơn vị.
\shortans{840502}
\loigiai{
\textbf{Đáp số:} $840502$\\
Số đó là $840\,502$.
}
\end{ex}

\begin{ex} %[4D3V10]
Có bao nhiêu số có 6 chữ số?
\shortans{900000}
\loigiai{
\textbf{Đáp số:} $900000$\\
$999\,999 - 100\,000 + 1 = 900\,000$.
}
\end{ex}

\begin{ex} %[4D3V10]
Tìm số chẵn bé nhất có 6 chữ số khác nhau.
\shortans{102346}
\loigiai{
\textbf{Đáp số:} $102346$\\
Số đó là $102\,346$.
}
\end{ex}

\begin{ex} %[4D3V10]
Tính tổng giá trị của chữ số 6 và chữ số 2 trong số $645\,208$.
\shortans{600200}
\loigiai{
\textbf{Đáp số:} $600200$\\
$600\,000 + 200 = 600\,200$.
}
\end{ex}

\begin{ex} %[4D3C10]
Tìm số tròn trăm nghìn có 6 chữ số, biết chữ số hàng trăm nghìn gấp 4 lần chữ số 2.
\shortans{800000}
\loigiai{
\textbf{Đáp số:} $800000$\\
Chữ số hàng trăm nghìn là $2 \times 4 = 8 \Rightarrow 800\,000$.
}
\end{ex}
